%% ── Section 3: DantinoX Architecture ────────────────────────────────────────
%%
%% Cite keys:
%%   ainslie-etal-2023-gqa  ·  austin-etal-2021-structured  ·
%%   bradbury-etal-2018-jax  ·  deepseekv2-2024  ·  fedus-etal-2022-switch  ·
%%   heek-etal-2024-flax  ·  hu-etal-2026-elf  ·  nie-etal-2024-llada  ·
%%   vaswani-etal-2017-attention  ·  hu-etal-2022-lora  ·  raffel-etal-2020-t5
%%   li-etal-2022-diffusion-lm  ·  lovelace-etal-2023-ld4lg  ·
%%   lee-etal-2026-flm  ·  chen-etal-2026-langflow  ·
%%   lou-etal-2023-discrete  ·  sahoo-etal-2024-mdlm
%% ─────────────────────────────────────────────────────────────────────────────

\section{DantinoX Architecture}

AR, discrete diffusion, and flow-matching models are typically developed in
separate codebases, conflating paradigm properties with implementation choices.
DantinoX eliminates this confound: a single shared backbone is parameterised by
attention type, feed-forward variant, positional encoding, normalisation, and
optional LoRA adapters; the generation paradigm is selected via a single
\texttt{paradigm} field, leaving all weights, hyperparameters, training loop,
and hardware layout unchanged.

\paragraph{Backbone.}
The shared backbone (Figure~\ref{fig:backbone}) is a pre-norm transformer with
$L$ identical blocks.
Each block contains a self-attention sublayer and a feed-forward sublayer
connected by residual paths.

Three attention mechanisms are available.
\textbf{Multi-Head Attention}
(MHA; \citealt{vaswani-etal-2017-attention}) serves as the reference
configuration.
\textbf{Grouped-Query Attention} (GQA; \citealt{ainslie-etal-2023-gqa})
partitions the $n_\mathrm{heads}$ query heads into groups that share a single
KV head, reducing the KV-cache by a factor of
$n_\mathrm{heads}/k_\mathrm{heads}$ at no cost to generation quality.
\textbf{Multi-Head Latent Attention} (MLA; \citealt{deepseekv2-2024}) stores a
low-dimensional latent vector per token and reconstructs keys and values
on-the-fly, reducing KV-cache size by a factor that depends on the chosen
latent dimension.
All three support Flash Attention, sliding-window context, and differential
attention as optional drop-in enhancements.

The feed-forward slot accepts either a dense MLP with SwiGLU or GELU activation
or a sparse Mixture-of-Experts layer \citep{fedus-etal-2022-switch} with top-$k$
routing.
Positional encoding is selected from RoPE, sinusoidal, learned, or none.
Normalisation is either RMSNorm or LayerNorm.
LoRA adapters \citep{hu-etal-2022-lora} can be attached to any projection
matrix in the backbone without modifying the shared weights.

\paragraph{Generation Paradigms.}
Each paradigm head connects to the backbone output
$\mathbf{h} \in \mathbb{R}^{B \times T \times D}$
(Figure~\ref{fig:paradigms}).
The three paradigms differ in attention masking, training objective, and
inference procedure; the backbone is unchanged.

\noindent\textbf{Autoregressive (AR).}
A causal mask restricts each position to its left context.
The training objective is next-token cross-entropy over all positions.
At inference, a static KV-cache confines each decoding step to a single forward
pass over the new token alone, holding additional memory at $\mathcal{O}(1)$
per step.
Greedy, top-$k$, and nucleus sampling are supported.

\noindent\textbf{Discrete Diffusion.}
DantinoX follows the LLaDA formulation \citep{nie-etal-2024-llada} for this
paradigm.
Building on prior masked diffusion language models
\citep{austin-etal-2021-structured,lou-etal-2023-discrete,sahoo-etal-2024-mdlm},
LLaDA is the first to scale masked diffusion to 8B parameters from scratch while
achieving performance competitive with autoregressive models, and it establishes the
$(1/t)$-weighted cross-entropy—an upper bound on the negative log-likelihood—as the
standard principled training objective.
The backbone runs with bidirectional attention.
During training, tokens are independently replaced by a \texttt{[MASK]} symbol
with probability $p_\mathrm{mask}(t)$ under a linear, cosine, or $\sqrt{\cdot}$
schedule \citep{austin-etal-2021-structured}.
The loss is a $(1/t)$-weighted cross-entropy over masked positions
\citep{nie-etal-2024-llada}.
At inference, $S$ reverse-diffusion passes unmask positions in order of
predicted confidence.
Four decoding strategies are provided: \texttt{sample}, \texttt{greedy},
\texttt{confidence}, and \texttt{factor}.
A block-generation mode denoises one fixed-width window at a time using a
DualCache that stores prefix and suffix KV states in separate buffers,
yielding a 1.4--2.1$\times$ throughput gain over global denoising
(Section~\ref{sec:results}).

\noindent\textbf{Continuous Flow-Matching.}
DantinoX follows the ELF formulation \citep{hu-etal-2026-elf} for this
paradigm.
Unlike prior continuous DLMs that apply per-step token-level cross-entropy
supervision along the denoising trajectory
\citep{li-etal-2022-diffusion-lm,nie-etal-2024-llada} or require a separate
decoder for final discretization \citep{lovelace-etal-2023-ld4lg}, and unlike
concurrent flow-based language models
\citep{lee-etal-2026-flm,chen-etal-2026-langflow}, ELF keeps the trajectory
entirely in continuous embedding space and discretizes only at the final step
via a shared-weight network.
The paradigm operates in the embedding space of a frozen T5 encoder
\citep{raffel-etal-2020-t5}.
The forward process is
$\mathbf{z}_t = t\,\mathbf{x} + (1{-}t)\,\boldsymbol\varepsilon$,
$t \sim \mathcal{U}[0,1]$,
where $\mathbf{x}$ is a normalised T5 embedding and
$\boldsymbol\varepsilon \sim \mathcal{N}(\mathbf{0},\mathbf{I})$.
The backbone learns the corresponding velocity field with a combined MSE
and cross-entropy loss \citep{hu-etal-2026-elf}.
Inference integrates the ODE by Euler steps; Classifier-Free Guidance is
controlled by a scalar \texttt{cfg\_scale} parameter.

\paragraph{Training.}
A \texttt{ModelConfig} object specifies the backbone architecture; wrapping it
in \texttt{Paradigm} selects the generation paradigm.
A \texttt{TrainingConfig} object specifies optimiser, schedule, batch size,
and hardware layout.
A single \texttt{Trainer(paradigm, tcfg).fit()} call launches training
(Listing~\ref{lst:training}).

\paragraph{Inference and Benchmarking.}
\texttt{Generator} reads the \texttt{ModelConfig} from a checkpoint directory
and dispatches to the appropriate algorithm—KV-cache decoding for AR,
reverse diffusion for discrete, and ODE integration for continuous flow-matching—without requiring
paradigm-specific code from the caller (Listing~\ref{lst:inference}).
\texttt{dx.count\_flops} computes an analytical per-component FLOPs breakdown
covering attention, FFN, and embedding layers without executing any GPU kernel.
\texttt{dx.profile} additionally measures wall-clock latency and Model FLOP
Utilisation (MFU) over configurable warm-up and measurement windows.
At the system level, \texttt{BenchmarkSuite.default()} sweeps throughput,
latency, and perplexity across batch sizes and sequence lengths and writes
structured CSV output for the analyses reported in Section~\ref{sec:results}.
